\documentclass[12pt]{article}
\usepackage[utf8]{inputenc}
\usepackage[T1]{fontenc}
\usepackage{hyperref}
\title{Sirva Language Resources for K-12 Curriculum Development}
\author{Compiled by Claude Sonnet 5.5 (claude-sonnet-5-5) for Andr\'as Kornai}
\date{}
\begin{document}
\maketitle

\section{Introduction}
Sirva (ISO 639-3: sbq; also called Sileibi) is a Papuan language of the Madang family (Southern Adelbert), spoken in Madang Province, Papua New Guinea, in villages including Kamambu, Kumbuna, Musita and Sileibi. Speaker figures are old (about 260 from 1981; about 259 in 1975). The main documentation is academic: Don Daniels' work on the Kamambu variety. No school or community-produced materials were found.

\section{Primary Educational Resources}
No Sirva primers, textbooks or school programmes were found. The language is not known to be taught in schools.

\section{Linguistic and Reference Materials}
\begin{itemize}
\item Daniels, Don (2015), Sirva grammar sketch, in \emph{A Reconstruction of Proto-Sogeram} (UC Santa Barbara dissertation), pp. 658--765; bibliographic records: \url{https://grambank.clld.org/sources/sDanielsSirva} and \url{https://endangeredlanguages.com/node/37278}
\item Daniels (2010), article on the phonological history of the Sogeram languages (ELP record): \url{https://endangeredlanguages.com/pt-br/node/37277}
\item New Guinea World profile of Sirva: \url{https://newguineaworld.linguistik.uzh.ch/families/madang-upper-yuat-river/madang/west-madang/sogeram-and-tomui-rivers/sogeram-river/mum-sirva/sirva}
\item Wikipedia, Sirva language: \url{https://en.wikipedia.org/wiki/Sirva_language}
\end{itemize}

\section{Religious and Community Resources}
Joshua Project reports no Bible portions, audio Bible or Jesus Film for the Sileibi people: \url{https://joshuaproject.net/people_groups/14176}

\section{Digital and Audio Resources}
PARADISEC holds archival items that appeared in searches (contents not verified):
\begin{itemize}
\item \url{https://catalog.paradisec.org.au/collections/DD1/items/579}
\item \url{https://catalog.paradisec.org.au/repository/DD1/343}
\end{itemize}
A University of Hawai`i ScholarSpace item also appeared: \url{https://scholarspace.manoa.hawaii.edu/items/2157d1a3-f2bb-44e8-a3e8-08e0411031a4}

\section{Curriculum Development Context}
With about 260 speakers, mostly adults, a curriculum would rely on Daniels' grammar sketch and archival recordings, and would need community-led development of an orthography and materials. Tok Pisin and English are the usual languages of schooling in PNG. Little exists beyond academic documentation.
\end{document}